\subsection{The undilated kernel and an explicit finite coefficient rule}
\label{mld:dil:undilated-algorithm}

The unequal-dilation theorem below uses a known undilated closure from
the incoming reports \cite{IncomingClosure,IncomingShifted}. We record its
generating function and derivation so that the resulting algorithm is
self-contained. These undilated statements are prior results within the
supplied material.

For $0<t<1$ and $\Re s,\Re v<1$, set
\[
 C(s,v;t)=\int_0^1\zeta(s,x)\zeta(v,\{x+t\})\dd x.
\]
The integral is ordinary: the two singularities are separated. Its two
equivalent continuations are
\begin{align}
 C(s,v;t)={}&\Gamma(1-s)\Gamma(1-v)(2\pi)^{s+v-2}
 \bigl\{e^{-\pi\ii(s-v)/2}\Li_{2-s-v}(e^{2\pi\ii t})\notag\\
 &\hspace{38mm}+e^{\pi\ii(s-v)/2}\Li_{2-s-v}(e^{-2\pi\ii t})\bigr\},
 \label{mld:dil:C-polylog}\\
 C(s,v;t)={}&\Gamma(s+v-1)
 \left\{\frac{\Gamma(1-s)}{\Gamma(v)}\zeta(s+v-1,t)
       +\frac{\Gamma(1-v)}{\Gamma(s)}\zeta(s+v-1,1-t)\right\}.
 \label{mld:dil:C-hurwitz}
\end{align}
All apparent singularities are interpreted in the complete expression.
To prove the first formula, start with $\Re s,\Re v<0$ and the absolutely
convergent Hurwitz Fourier series, whose nonzero coefficients are
\[
 \widehat\zeta_s(k)=\Gamma(1-s)(2\pi\ii k)^{s-1},\qquad k\ne0.
\]
The zero coefficient is zero. Integrating the product retains the
opposite frequencies and yields \eqref{mld:dil:C-polylog}. For the second
formula, insert the same Fourier expansion with spectral order $s+v-1$
on its right. Euler's gamma reflection identity reduces the coefficient
of the positive-argument polylogarithm by
\[
 \frac{\sin(\pi s)e^{-\pi\ii(s+v)/2}
       +\sin(\pi v)e^{\pi\ii(s+v)/2}}{\sin\pi(s+v)}
       =e^{-\pi\ii(s-v)/2}.
\]
The other coefficient is obtained by reversing the signs in the phases.
This calculation holds off the integer hyperplanes; analytic continuation
extends it across their removable singularities. The local decomposition
$\zeta(s,x)=x^{-s}+\zeta(s,1+x)$ justifies joint holomorphy of the
integral on $\Re s,\Re v<1$, and hence the continuation to that region.
This is the classical Fourier and product-integral mechanism
\cite{DLMFhurwitz,EspinosaMoll} in the circular-translation convention.

Define, near $(u,v)=(0,0)$,
\begin{align*}
 \mathcal A(u,v)&=
  \frac{\Gamma(1+u+v)\Gamma(1-v)}{\Gamma(1+u)},
 &\mathcal B(u,v)&=\mathcal A(v,u),\\
 \mathcal K(u,v;t)&=uv\,C(1+u,1+v;t).
\end{align*}
Equation~\eqref{mld:dil:C-hurwitz} gives
\begin{equation}
 \mathcal K(u,v;t)=
 -u\mathcal A(u,v)\zeta(1+u+v,1-t)
 -v\mathcal B(u,v)\zeta(1+u+v,t).
 \label{mld:dil:K-hurwitz}
\end{equation}
This function is holomorphic at the origin. Indeed,
$\mathcal A(u,-u)=\mathcal B(u,-u)=1$, so
$(u\mathcal A+v\mathcal B)/(u+v)$ has a removable denominator. Removing
the common Hurwitz pole in \eqref{mld:dil:K-hurwitz} leaves the negative of that
quotient plus analytic Stieltjes series.

For the canonical undilated finite part one consequently has the
terminating coefficient rule
\begin{equation}
 \boxed{I_{mn}(t)=(-1)^{m+n}m!n!
       [u^{m+1}v^{n+1}]\mathcal K(u,v;t).}
 \label{mld:dil:I-coefficient}
\end{equation}
Here $I_{mn}$ denotes the constant term of
$\int\gamma_m(x)\gamma_n(\{x+t\})\dd x$ with the two ambient cutoffs.
For clarity, the matching to a finite part can be justified before
extracting any coefficients. Put
$\mathcal H(u,x)=\zeta(1+u,x)-1/u$ and $b=1-t$.
The following locally convergent expression is analytic for small $u,v$:
\begin{align*}
 &\int_0^b\left\{
 x^{-1-u}[\mathcal H(v,t+x)-\mathcal H(v,t)]
       +\mathcal H(u,1+x)\mathcal H(v,t+x)\right\}\dd x\\
 &+\int_0^t\left\{
 y^{-1-v}[\mathcal H(u,b+y)-\mathcal H(u,b)]
       +\mathcal H(v,1+y)\mathcal H(u,b+y)\right\}\dd y\\
 &+\mathcal H(v,t)\frac{1-b^{-u}}u
       +\mathcal H(u,b)\frac{1-t^{-v}}v.
\end{align*}
The bracketed differences vanish linearly; both quotients have removable
values at zero. The coefficient of $u^mv^n$ is exactly
$(-1)^{m+n}I_{mn}(t)/(m!n!)$, by direct integration of the subtracted
logarithmic singularities. In a left spectral half-plane the displayed
expression is also
\[
 C(1+u,1+v;t)+\frac1{uv}
       +\frac{\mathcal H(v,t)}u+\frac{\mathcal H(u,b)}v.
\]
Using the axis values of \eqref{mld:dil:K-hurwitz}, it equals
\[
 \frac{\mathcal K(u,v;t)-\mathcal K(u,0;t)
            -\mathcal K(0,v;t)+\mathcal K(0,0;t)}{uv}.
\]
This proves \eqref{mld:dil:I-coefficient} with the specified normalization.

The coefficient calculation uses rational polynomial arithmetic and
positive integer zeta values only, because
\begin{equation}
 \log\mathcal A(u,v)=\sum_{j\ge2}\frac{\zeta(j)}j
       \left\{(-1)^j\bigl((u+v)^j-u^j\bigr)+v^j\right\}.
 \label{mld:dil:logA}
\end{equation}
For an explicit finite algorithm, let $\ell_j$ be the homogeneous summand
of degree $j$ here. The homogeneous terms of $\mathcal A$ satisfy
\[
 \mathcal A_0=1,\qquad\mathcal A_1=0,\qquad
 \mathcal A_d=\frac1d\sum_{j=2}^d j\ell_j\mathcal A_{d-j}.
\]
Only $d\le m+n+2$ can contribute to \eqref{mld:dil:I-coefficient}. This proves
finite linear closure in the single values $\gamma_j(t)$ and
$\gamma_j(1-t)$, with $j\le m+n+1$. Their top coefficients are
$1/(m+1)$ and $1/(n+1)$; the coefficients at order $m+n$ vanish since
$\mathcal A_1=0$. These are the incoming closure theorem and its
missing-order rule, now available as an explicit input to the new
dilation calculation.
